\section{A nonconstructive family with a different first factor}
\label{sec:nonconstructive}
The graphs in Theorems~\ref{thm:C} and~\ref{thm:L} are determined before
their maximality is proved. We next give a different construction in which
the first factor is selected by Zorn's lemma. Its premises are verified
by a countable set of specified rows. An additional row ensures that the
selected factor differs from the corresponding determined factor, for
every allowed selection. The endpoint obstruction remains the same;
this is not an independent mechanism for the sum question.

\begin{lemma}[An extension preserving a lower pairing bound]
\label{lem:capped-extension}
Let $X$ be a real Banach space and let $S\subset X\times X^*$ be a
nonempty monotone relation. Suppose $c,\rho>0$ satisfy
\[
\pair{x}{a}\ge-c\quad((x,a)\in S),\qquad
S^\mu\cap\bigl(\overline B_X(0,\rho)\times X^*\bigr)=\varnothing.
\]
There is a maximally monotone graph $G\supset S$ such that
\[
\pair{x}{a}\ge-c,\qquad \|x\|>\rho\quad((x,a)\in G).
\]
In particular, $0\notin\dom G$ and $\mathcal V(G)\le c/\rho$.
\end{lemma}
\begin{proof}
Order by inclusion the monotone relations containing $S$ whose every row
satisfies the lower pairing bound $-c$. The relation $S$ belongs to this
family. The union of a chain is monotone, since any two of its rows occur
together in one chain member; it also contains $S$ and preserves the
bound. Zorn's lemma supplies an inclusion-maximal member $G$.

We prove maximality in the full pair $X\times X^*$, not merely in this
ordered family. For pairs $u=(x,a)$ and $v=(y,b)$, write
$[u,v]=\pair{x-y}{a-b}$, and let $o=(0,0)$. If
$w\in G^\mu\setminus G$ and $[o,w]\ge-c$, adjoining $w$ contradicts
the choice of $G$. Otherwise put $d=-[o,w]>c$ and $\alpha=c/d\in(0,1)$.
For every $v\in G$, expansion of the pairing gives
\begin{align*}
[\alpha w,v]
&=(1-\alpha)[o,v]+\alpha[w,v]-\alpha(1-\alpha)[o,w]\\
&\ge-(1-\alpha)c+\alpha(1-\alpha)d=0.
\end{align*}
Thus $\alpha w\in G^\mu$, and
$[o,\alpha w]=-c^2/d>-c$. But
$[w,\alpha w]=-(1-\alpha)^2d<0$, so $\alpha w\notin G$ because
$w$ is related to every point of $G$. Adjoining $\alpha w$ is another
contradiction. Therefore $G^\mu=G$.

Since $S\subset G$, polar reversal gives $G=G^\mu\subset S^\mu$.
The assumed exclusion implies $\|x\|>\rho$ on every row. The lower
pairing bound then gives
$\max\{0,-\pair{x}{a}\}/\|x\|\le c/\rho$. Taking the supremum
proves the last assertion.
\end{proof}

\begin{lemma}[A full-domain norm partner]\label{lem:norm-partner}
Let $M:X\rightrightarrows X^*$ be maximally monotone, with
$0\notin\dom M$, on a real Banach space. For $\lambda>0$, let
$N_\lambda=\partial(\lambda\|\cdot\|)$. This operator is maximally
monotone and has full domain. If $\mathcal V(\gra M)\le\lambda$, then
$(M,N_\lambda)$ satisfies the interior-domain condition and
$M+N_\lambda$ is not maximally monotone, with polar witness $(0,0)$.
\end{lemma}
\begin{proof}
The subgradient inequality, tested at zero and at scalar multiples of
$x$, gives $\pair{x}{j}=\lambda\|x\|$. Substitution back into that
inequality gives $\pair{y}{j}\le\lambda\|y\|$ for every $y\in X$,
which is equivalent to $\|j\|\le\lambda$. Thus
\[
N_\lambda x=\{j\in X^*: \|j\|\le\lambda,
                  \pair{x}{j}=\lambda\|x\|\}.
\]
Conversely, these two conditions imply the subgradient inequality by the
dual norm bound. Hahn--Banach provides a norming functional for every
$x\ne0$, and $N_\lambda0=\lambda\overline B_{X^*}$. Hence the domain is
all of $X$. For $j\in N_\lambda x$ and $k\in N_\lambda y$,
\[
\pair{x-y}{j-k}
=\lambda\|x\|+\lambda\|y\|-\pair{x}{k}-\pair{y}{j}\ge0,
\]
so the graph is monotone.

Suppose $(z,p)$ is related to this entire graph. For any $v\in X$,
choose $j\in N_\lambda(z+v)$. The test at $z+v$ gives
$\pair{v}{p}\le\pair{v}{j}\le\lambda\|v\|$. Replacing $v$ by $-v$
shows $\|p\|\le\lambda$. If $z\ne0$, testing at $tz$ for $t>0$ yields
\[
(1-t)\bigl(\pair{z}{p}-\lambda\|z\|\bigr)\ge0.
\]
Values on both sides of $t=1$ force equality. At $z=0$ that equality is
automatic. Thus $p\in N_\lambda z$, proving maximality directly.
This is the norm case of the subdifferential maximality theorem
\cite{rockafellar1970subdiff}.

For every $(x,a)\in\gra M$ and every $j\in N_\lambda x$,
\[
\pair{x}{a+j}=\pair{x}{a}+\lambda\|x\|\ge0.
\]
Thus $(0,0)$ is related to every sum row. It is absent from the graph
because $M0=\varnothing$. The sum is monotone, and adjoining that point
gives a proper monotone extension. Finally, a maximally monotone graph is
nonempty, since the empty graph admits a one-point extension. Together
with $\dom N_\lambda=X$, this proves the ordered intersection condition.
\end{proof}

We now use the concrete maps of Eqs.~\eqref{eq:D4}--\eqref{eq:D9}.
Let $K_{\mathbb Q}$ be the set of finitely supported rational labels
$k\in\ell^1(I)$ satisfying $mk=0$ and $r(k)=0$. Put
\[
t_n^+=1+1/n,\quad t_n^-=-1-1/n\quad(n\ge1),\qquad
S_0=\{(La,a):a=a^{t_n^\pm}+k,\ n\ge1,\ k\in K_{\mathbb Q}\}.
\]
Define one more label and the augmented relation by
\begin{equation}
p_\circ=-4e_{0,1}+4e_{0,3},\qquad x_\circ=Lp_\circ,
\qquad S=S_0\cup\{(x_\circ,p_\circ)\}.
\label{eq:nonconstructive-seed}
\end{equation}

\renewcommand{\thetheorem}{N}
\begin{theorem}[A selected counterexample on both spaces]
\label{thm:nonconstructive}
The relation $S$ is a countable monotone subset of
$c_0(I)\times\ell^1(I)$. There exists a maximally monotone
$\mathcal A$ whose graph contains $S$ and whose every row satisfies
\[
\|x\|_\infty>1/2,\qquad \pair{x}{a}\ge-\sigma,
\qquad \mathcal V(\gra\mathcal A)\le2\sigma.
\]
Every such extension satisfies $\mathcal A\ne A$, and
$\mathcal A+\partial(2\sigma\|\cdot\|_\infty)$ is nonmaximal under
the original interior-domain condition.
With the specified $Q:\ell^1\to c_0(I)$, the operator
$\mathcal T=Q^*\mathcal A Q$ is maximally monotone in
$\ell^1\times\ell^\infty$, differs from $T$, and satisfies
\[
0\notin\dom\mathcal T,\qquad
\mathcal V(\gra\mathcal T)\le2\sigma.
\]
The pair $(\mathcal T,\partial(2\sigma\|\cdot\|_1))$ also satisfies
the original condition and has a nonmaximal sum. Both sums have
$(0,0)$ in their full monotone polar and outside their graph.
\end{theorem}
\begin{proof}
Finite rational labels form a countable set. Each $a^{t_n^\pm}$ has
finite support, so $S_0$ is countable. It is a subset of $\gra A$,
because adding $k\in K_{\mathbb Q}\subset K$ leaves $r$ and $m$
unchanged. Thus $S_0$ is nonempty and monotone. Its every row has
pairing greater than $-\sigma$ by Eq.~\eqref{eq:C9}.
The added label satisfies
\[
r(p_\circ)=4,\qquad mp_\circ=0,\qquad
x_\circ=-8e_{0,1}-12e_{0,2}-4e_{0,3},\qquad
\pair{x_\circ}{p_\circ}=16.
\]
For a row $(La,a)\in S_0$, put $t=r(a)=t_n^\pm$. Since $|t|\le2$,
Eq.~\eqref{eq:C2}, valid for all dual labels, gives
\[
\pair{x_\circ-La}{p_\circ-a}
=(4-t)^2-\|F(t)\|_2^2\ge3-\sigma>0.
\]
Consequently $S$ is monotone and obeys the lower pairing bound
$-\sigma$. Also $(x_\circ,p_\circ)\notin\gra A$, since
$mp_\circ=0\ne F(4)=(1,0,\ldots)$.

It remains to prove the full-polar exclusion required by
Lemma~\ref{lem:capped-extension}. Take any
$(x,p)\in S_0^\mu$, with $x\in c_0(I)$ and $p\in\ell^1(I)$.
Fixing any $t_n^\pm$ and testing the labels
$a^{t_n^\pm}+j k$ for all integers $j$ and all $k\in K_{\mathbb Q}$
in Eq.~\eqref{eq:C4} forces $\pair{x+Ep}{k}=0$.
All coordinate differences used in the annihilator calculation in the
proof of Theorem~\ref{thm:C} belong to $K_{\mathbb Q}$: they are
$e_{b,j}-e_{b,k}$ for $b\ge1$, those with $j,k\ge3$ in block zero,
and $g=e_{0,1}-e_{0,2}$. Their tests and the global $c_0$ condition
therefore give $x=-Ep+v h$ for some $v\in\R$.
In particular, $\pair{x}{g}=r(p)$ by Eq.~\eqref{eq:C1} and $r(g)=0$.

Suppose $\rho:=r(p)$ satisfies $|\rho|\le1$, and write $mp=(b,z)$.
For the actual tests $t=t_n^\pm$, Eq.~\eqref{eq:C4} becomes
\[
v\rho-(v+\rho)t+t^2-\|F(t)-mp\|_2^2\ge0.
\]
Use $\omega_{1/n}\to\omega_0$ in $\ell^2$, and put
$\eta=\omega_0$. Letting $n\to\infty$ on the two branches gives
\begin{align*}
v(\rho-1)-\rho+2b-b^2-\|z-\eta\|_2^2&\ge0,\\
v(\rho+1)+\rho-2b-b^2-\|z-\eta\|_2^2&\ge0.
\end{align*}
Multiply by $(1+\rho)/2$ and $(1-\rho)/2$, respectively, and add.
These are nonnegative weights even at $\rho=\pm1$; the terms in $v$
cancel, leaving
\[
-(b-\rho)^2-\|z-\eta\|_2^2\ge0.
\]
This forces $z=\eta$, which is impossible because $mp\in\ell^1$
whereas $\eta\notin\ell^1$. Thus every point of $S_0^\mu$ has
$|\pair{x}{g}|>1$ and $\|x\|_\infty>1/2$, since $\|g\|_1=2$.
Polar reversal gives the same exclusion for $S^\mu\subset S_0^\mu$.
All limits were comparisons with actual seed rows, not additions of
endpoint labels to the dual space.

Apply Lemma~\ref{lem:capped-extension} with $c=\sigma$ and radius
$1/2$. It produces $\mathcal A$ with the asserted maximality and
bounds. The forced row $(x_\circ,p_\circ)$ proves
$\mathcal A\ne A$ for every such choice. Lemma~\ref{lem:norm-partner}
proves the claimed counterexample on $c_0$. More explicitly, for every
$(x,a)\in\gra\mathcal A$ and every norm-subgradient value $j$,
\[
\pair{x}{a+j}\ge-\sigma+2\sigma\|x\|_\infty>0.
\]

For completeness, define the entire pullback graph by
\[
\gra\mathcal T=\{(u,Q^*a):(Qu,a)\in\gra\mathcal A\}.
\]
The surjectivity and norm-$2$ lifting bound of $Q$ were proved in
Eqs.~\eqref{eq:L1}--\eqref{eq:L2}, and the full-dual adjoint isometry
in Eq.~\eqref{eq:L3}; they do not depend on which monotone graph is
pulled back. This graph is nonempty and monotone by the adjoint identity.
For any polar query $(v,v^*)\in\ell^1\times\ell^\infty$, all translates
$u+t k$, $k\in\ker Q$, of one actual graph row force
$\pair{k}{v^*}=0$. The functional
$x\mapsto\pair{u}{v^*}$, $Qu=x$, is therefore well-defined and linear.
Choosing a norm-$2$ lift bounds it by $2\|v^*\|_\infty\|x\|_\infty$.
The full dual of $c_0(I)$ supplies $p\in\ell^1(I)$ representing this
functional, so $v^*=Q^*p$. Every row of $\mathcal A$ has a lift, and
the remaining polar inequalities are
\[
\pair{Qv-x}{p-a}\ge0\qquad((x,a)\in\gra\mathcal A).
\]
Maximality of $\mathcal A$ gives $(Qv,p)\in\gra\mathcal A$ and hence
$(v,v^*)\in\gra\mathcal T$. This proves full-dual maximality.

For every pullback row, $\|u\|_1\ge\|Qu\|_\infty>1/2$ and
$\pair{u}{Q^*a}=\pair{Qu}{a}\ge-\sigma$. Thus the asserted missing
origin and radial bound hold. A lift $u_\circ$ of $x_\circ$ gives a row
$(u_\circ,Q^*p_\circ)$ of $\mathcal T$. If it were a row of $T$,
injectivity of $Q^*$ would force its original label to be $p_\circ$,
contradicting $p_\circ\notin D$. Therefore $\mathcal T\ne T$.
Lemma~\ref{lem:norm-partner} now proves the second sum claim, including
full domain of its norm partner, all graph values, the ordered
intersection condition, and graph nonmembership of $(0,0)$.
\end{proof}

This theorem selects an unspecified member of a nonempty family. We do
not claim that two selections give different operators or that membership
in the selected graph is decidable. The forced row proves exactly the
distinction stated in the theorem. The proof also explains why a lower
pairing bound on seed rows alone is insufficient: the radial argument
is needed to promote a bound-preserving extension to global maximality.
